\itemhead{复习题二 · 第3题（教材p.61）}{XJ-B1-C02-RV-P061-Y0779}
\label{XJ-B1-C02-RV-P061-Y0779}
{\footnotesize 湘教版教材p.61 / PDF第68页 / 复习题二 / 第3题}
\field{原题}{证明下列不等式：
\subq{1}若 $a,b,c,d$ 都是正数，则 $(ab+cd)(ac+bd)\ge4abcd$；
\subq{2}若 $a,b,c$ 是非负实数，则 $a(b^2+c^2)+b(c^2+a^2)+c(a^2+b^2)\ge6abc$；
\subq{3}若 $a,b$ 是非负实数，则 $a+b+2\ge2(\sqrt a+\sqrt b)$；
\subq{4}若 $a,b\in\mathbb R$，则 $\dfrac{a^2+b^2}2\ge\left(\dfrac{a+b}2\right)^2$。}
\field{参考答案}{（1）分别对两因式用基本不等式得 $\ge(2\sqrt{abcd})^2$，等号 $a=d,b=c$；（2）将左边减去 $6abc$ 写为 $a(b-c)^2+b(c-a)^2+c(a-b)^2\ge0$；（3）差为 $(\sqrt a-1)^2+(\sqrt b-1)^2\ge0$，等号 $a=b=1$；（4）差为 $(a-b)^2/4\ge0$，等号 $a=b$。}
{\small \field{关联知识}{I03 基本不等式、等号与最值建模（教材帮p.52、53、54 / PDF61、62、63）}
\field{方法与易错点}{列正数条件与固定量，构造基本不等式，解等号条件并代回实际限制。 易错：把一个下界误当最小值；多个等号条件不能同时实现；漏单位和整数限制。}
\field{答案依据}{答案册相关题组：PDF第15、16页（册内页14、15）；按当前题库原题逐项重算，不按旧版题号机械对套。}}
